
\chapter{Processing of Telos Frames: libtelos}
\label{cha:libtelos}

This chapter explains the library {\tt libtelos},
which contains the Telos parser. The Telos parser is able to parse the answers in FRAME
or LABEL format from ConceptBase.

To call the Telos parser, you must link your program with the library
{\tt libtelos.a/libtelos.dll} which can be found in the directory
\verb+$CB_HOME/<arch>/lib+ where \verb+<arch>+ is either sun4, i86pc,
linux, or windows.

In your source files, you must include the header files
{\tt fragment.h}, {\ttfamily te\_access.h}, {\tt te\_callparser.h},
{\tt te\_cursor.h}, and/or {\tt te\_smlutil.h}.
All header files are located in the directory \verb+$CB_HOME/include+.

The following sections explain several functions to call the
parser and to handle the data structures. In principle, there are three
different ways to parse and to access Telos frames:
\begin{itemize}
\item Using the functions and data structures defined in {\tt fragment.h},
{\tt te\_callparser.h}, and {\tt te\_smlutil.h}: The Telos parser is invoked
directly and the contents of the Telos frames is retrieved by navigating
over a list of fragments (a fragment is a data structure for a Telos frame).
See section \ref{sec:fragmenth} for details.
\item Using the functions and data structures defined in {\ttfamily te\_access.h}:
Telos frames are represented in vectorized structure. One can use functions
to create, destroy or apply to filters to the structure. See section \ref{sec:accessh} for details.
\item Using the functions and data structures defined in {\ttfamily te\_cursor.h}:
Iterating over a set of Telos frames is done by using a cursor.
See section \ref{sec:cursorh} for details.
\end{itemize}

The documentation in the following sections has been generated with
DOC++ (\url{http://docpp.sourceforge.net}).
% man kann mit CB_Make docpptex ein TeX-File generieren, dass die Dokumentation enth�lt.
% Dieses File wurde aber nachtraeglich editiert, daher sollte man Aenderungen besser
% manuell hier nachtragen.

\section{fragment.h and te\_callparser.h}
\label{sec:fragmenth}

\begin{cxxentry}
{typedef\ struct\ }
        {bindingList}
        {BindingList}
        {A binding list represents the list of parameters in a derive expression}
        {1}
\begin{cxxdoc}
A binding list represents the list of parameters in a
derive expression
\end{cxxdoc}
\end{cxxentry}
\begin{cxxentry}
{typedef\ struct\ }
        {objectIdentifier}
        {ObjectIdentifier}
        {An object identifier represents a Telos object name.}
        {2}
\begin{cxxdoc}
An object identifier represents a Telos object name.
It may be a simple object name, a derive expression,
or a select expression.
\end{cxxdoc}
\end{cxxentry}
\begin{cxxentry}
{typedef\ struct\ }
        {classlist}
        {te\_ClassList}
        {A class list is a list of object identifiers }
        {3}
\begin{cxxdoc}
A class list is a list of object identifiers
\end{cxxdoc}
\end{cxxentry}
\begin{cxxentry}
{typedef\ struct\ }
        {attrclasslist}
        {AttrClassList}
        {An AttrClassList is a list of attribute categories.}
        {4}
\begin{cxxdoc}
An AttrClassList is a list of attribute categories.
Attribute categories or simple labels.
\end{cxxdoc}
\end{cxxentry}
\begin{cxxentry}
{typedef\ struct\ }
        {specObjId}
        {SpecObjId}
        {Used only internally for extended syntax }
        {5}
\begin{cxxdoc}
Used only internally for extended syntax
\end{cxxdoc}
\end{cxxentry}
\begin{cxxentry}
{typedef\ struct\ }
        {selectexpb}
        {SelectExpB}
        {Used only internally for extended syntax }
        {6}
\begin{cxxdoc}
Used only internally for extended syntax
\end{cxxdoc}
\end{cxxentry}
\begin{cxxentry}
{typedef\ struct\ }
        {restriction}
        {Restriction}
        {Used only internally for extended syntax }
        {7}
\begin{cxxdoc}
Used only internally for extended syntax
\end{cxxdoc}
\end{cxxentry}
\begin{cxxentry}
{typedef\ struct\ }
        {objectset}
        {ObjectSet}
        {Used only internally for extended syntax }
        {8}
\begin{cxxdoc}
Used only internally for extended syntax
\end{cxxdoc}
\end{cxxentry}
\begin{cxxentry}
{typedef\ struct\ }
        {propertylist}
        {PropertyList}
        {A property list is a list of attributes.}
        {9}
\begin{cxxdoc}
A property list is a list of attributes. Attributes
have a label and a value. The member objectSet is used
only in an extended syntax.
\end{cxxdoc}
\end{cxxentry}
\begin{cxxentry}
{typedef\ struct\ }
        {attrdecllist}
        {AttrDeclList}
        {An AttrDeclList is a list of attribute declarations.}
        {10}
\begin{cxxdoc}
An AttrDeclList is a list of attribute declarations.
It represents everthing between "with" and "end" in a
Telos frame. One attribute declaration has a list of attribute
categories and a list of properties (attribute definitions).
\end{cxxdoc}
\end{cxxentry}
\begin{cxxentry}
{typedef\ struct\ }
        {smlfragmentList}
        {te\_SMLfragmentList}
        {A SMLfragmentList is a list of Telos frames.}
        {11}
\begin{cxxdoc}
A SMLfragmentList is a list of Telos frames. Each Telos
frame has an object identifier (id). It may have in addition
an inOmega class, a list of in-Classes, a list of isA-Classes,
and an attribute declaration. Except id, all members may be NULL.
\end{cxxdoc}
\end{cxxentry}
\begin{cxxentry}
{typedef\ struct\ }
        {frameParseoutput}
        {FrameParseOutput}
        { FrameParseOutput is the structure returned by the function te\_frame\_parser.}
        {12}
\begin{cxxnames}
\cxxitem{te\_SMLfragmentList*}
        {smlfrag}
        {}
        {the list of fragments }
        {12.1}
\cxxitem{int}
        {error}
        {}
        {0 if ok, 1 if parse error, 2 if input is null }
        {12.2}
\cxxitem{char*}
        {errortoken}
        {}
        {If there was an parse error, this should indicate the token that caused the error.}
        {12.3}
\cxxitem{int}
        {errorline}
        {}
        {If there was an parse error, this should be the line number of the error.}
        {12.4}
\end{cxxnames}
\begin{cxxdoc}

FrameParseOutput is the structure returned by the function
te\_frame\_parser. It contains either a list of fragments or
information about the parse error.

\end{cxxdoc}
\begin{cxxvariable}
{te\_SMLfragmentList*}
        {smlfrag}
        {}
        {the list of fragments }
        {12.1}
\end{cxxvariable}
\begin{cxxvariable}
{int}
        {error}
        {}
        {0 if ok, 1 if parse error, 2 if input is null }
        {12.2}
\end{cxxvariable}
\begin{cxxvariable}
{char*}
        {errortoken}
        {}
        {If there was an parse error, this should indicate the token that caused the error.}
        {12.3}
\end{cxxvariable}
\begin{cxxvariable}
{int}
        {errorline}
        {}
        {If there was an parse error, this should be the line number of the error.}
        {12.4}
\end{cxxvariable}
\end{cxxentry}
\begin{cxxentry}
{typedef\ struct\ }
        {classlistParseoutput}
        {ClassListParseOutput}
        { ClassListParseOutput is the structure returned by the function te\_classlist\_parser.}
        {13}
\begin{cxxnames}
\cxxitem{te\_ClassList*}
        {classlist}
        {}
        {A list of classes (object names) }
        {13.1}
\cxxitem{int}
        {error}
        {}
        {Non-zero if an error occured }
        {13.2}
\cxxitem{char*}
        {errortoken}
        {}
        {If there was an parse error, this should indicate the token that caused the error.}
        {13.3}
\cxxitem{int}
        {errorline}
        {}
        {If there was an parse error, this should be the line number of the error.}
        {13.4}
\end{cxxnames}
\begin{cxxdoc}

ClassListParseOutput is the structure returned by the function
te\_classlist\_parser. It contains either a list of classes or
information about the parse error.
\end{cxxdoc}
\begin{cxxvariable}
{te\_ClassList*}
        {classlist}
        {}
        {A list of classes (object names) }
        {13.1}
\end{cxxvariable}
\begin{cxxvariable}
{int}
        {error}
        {}
        {Non-zero if an error occured }
        {13.2}
\end{cxxvariable}
\begin{cxxvariable}
{char*}
        {errortoken}
        {}
        {If there was an parse error, this should indicate the token that caused the error.}
        {13.3}
\end{cxxvariable}
\begin{cxxvariable}
{int}
        {errorline}
        {}
        {If there was an parse error, this should be the line number of the error.}
        {13.4}
\end{cxxvariable}
\end{cxxentry}
\begin{cxxfunction}
{FrameParseOutput*}
        {te\_frame\_parser}
        {(char*\ indata)}
        {Calls the Telos Parser to parse frames.}
        {14}
\cxxParameter{
\begin{tabular}[t]{lp{0.5\textwidth}}
{\tt\strut indata} & a string containing the input frames
\end{tabular}}
\cxxReturn{
\begin{tabular}[t]{lp{0.5\textwidth}}
{\tt\strut a} & pointer to a FrameParseOut structure
\end{tabular}}
\begin{cxxdoc}
Calls the Telos Parser to parse frames.

\end{cxxdoc}
\end{cxxfunction}
\begin{cxxfunction}
{ClassListParseOutput*}
        {te\_classlist\_parser}
        {(char*\ indata)}
        {Calls the Telos Parser to parse a list of object names.}
        {15}
\cxxParameter{
\begin{tabular}[t]{lp{0.5\textwidth}}
{\tt\strut indata} & a string containing the object names
\end{tabular}}
\cxxReturn{
\begin{tabular}[t]{lp{0.5\textwidth}}
{\tt\strut a} & pointer to a ClassListParseOut structure
\end{tabular}}
\begin{cxxdoc}
Calls the Telos Parser to parse a list of object names.

\end{cxxdoc}
\end{cxxfunction}
\begin{cxxfunction}
{char*}
        {FragmentToString}
        {(te\_SMLfragmentList *cursor}
        {Unparse a fragment list into a string}
        {16}
\begin{cxxdoc}

Unparse a fragment list into a string
\end{cxxdoc}
\end{cxxfunction}
\begin{cxxfunction}
{void}
        {DestroySMLfrag}
        {(te\_SMLfragmentList*\ fragment)}
        { Destroy a fragment list}
        {17}
\begin{cxxdoc}

Destroy a fragment list
\end{cxxdoc}
\end{cxxfunction}
\begin{cxxfunction}
{void}
        {Destroy\_ClassList}
        {(te\_ClassList*\ clist)}
        { Destroy a class list}
        {18}
\begin{cxxdoc}

Destroy a class list
\end{cxxdoc}
\end{cxxfunction}



\section{te\_access.h}
\label{sec:accessh}


\begin{cxxunion}
{struct\ }
        {te\_AttrDecl}
        {}
        {This structure represents an attribute declaration.}
        {19}
\begin{cxxnames}
\cxxitem{char**}
        {aszCategory}
        {}
        {contains the list of category labels }
        {19.1}
\cxxitem{char**}
        {aszLabel}
        {}
        {contains the list of property labels corresponding to }
        {19.2}
\cxxitem{char**}
        {aszValue}
        {}
        {the list of property values }
        {19.3}
\end{cxxnames}
\begin{cxxdoc}
This structure represents an attribute declaration. An attribute declaration
is a list of attribute categories with a list of properties (label and values)
that belong to these attribute categories.
\end{cxxdoc}
\begin{cxxvariable}
{char**}
        {aszCategory}
        {}
        {contains the list of category labels }
        {19.1}
\end{cxxvariable}
\begin{cxxvariable}
{char**}
        {aszLabel}
        {}
        {contains the list of property labels corresponding to }
        {19.2}
\end{cxxvariable}
\begin{cxxvariable}
{char**}
        {aszValue}
        {}
        {the list of property values }
        {19.3}
\end{cxxvariable}
\end{cxxunion}
\begin{cxxentry}
{typedef\ struct\ }
        {te\_AttrDecl}
        {TAttrDecl}
        {The type for te\_AttrDecl }
        {20}
\end{cxxentry}
\begin{cxxentry}
{typedef\ }
        {TAttrDecl*}
        {PAttrDecl}
        {The pointer for TAttrDecl }
        {21}
\begin{cxxdoc}
The pointer for TAttrDecl
\end{cxxdoc}
\end{cxxentry}
\begin{cxxentry}
{typedef\ struct\ }
        {te\_VectorizedTelosframe}
        {VTelos}
        {The type for te\_VectorizedTelosframe }
        {22}
\begin{cxxdoc}
The type for te\_VectorizedTelosframe
\end{cxxdoc}
\end{cxxentry}
\begin{cxxentry}
{typedef\ }
        {VTelos*}
        {PVTelos}
        {A pointer to VTelos }
        {23}
\begin{cxxdoc}
A pointer to VTelos
\end{cxxdoc}
\end{cxxentry}
\begin{cxxentry}
{typedef\ }
        {PVTelos*}
        {AVTelos}
        {An array of VTelos pointers }
        {24}
\begin{cxxdoc}
An array of VTelos pointers
\end{cxxdoc}
\end{cxxentry}
\begin{cxxunion}
{struct\ }
        {te\_TelosReport}
        {}
        {A te\_TelosReport is a projection on certain attributes of a frame }
        {25}
\begin{cxxnames}
\cxxitem{char**}
        {aszLabel}
        {}
        {List of labels in the report }
        {25.1}
\cxxitem{char**}
        {aszValue}
        {}
        {List of values in the report }
        {25.2}
\end{cxxnames}
\begin{cxxdoc}
A te\_TelosReport is a projection on certain attributes of a frame
\end{cxxdoc}
\begin{cxxvariable}
{char**}
        {aszLabel}
        {}
        {List of labels in the report }
        {25.1}
\end{cxxvariable}
\begin{cxxvariable}
{char**}
        {aszValue}
        {}
        {List of values in the report }
        {25.2}
\end{cxxvariable}
\end{cxxunion}
\begin{cxxentry}
{typedef\ struct\ }
        {te\_TelosReport}
        {TReport}
        {The type for te\_TelosReport }
        {26}
\end{cxxentry}
\begin{cxxentry}
{typedef\ }
        {TReport*}
        {PReport}
        {A pointer to TReport}
        {27}
\end{cxxentry}
\begin{cxxfunction}
{AVTelos}
        {vt\_createByFragment}
        {(te\_SMLfragmentList*\ fl)}
        { Maps the given fragmentlist fl into a vector of frames.}
        {28}
\cxxParameter{
\begin{tabular}[t]{lp{0.5\textwidth}}
{\tt\strut fl} & contains the fragment list in the way produced by the parser
\end{tabular}}
\cxxReturn{
\begin{tabular}[t]{lp{0.5\textwidth}}
{\tt\strut NULL} & if the argument is NULL too, else the pointer to the vector tree structure
\end{tabular}}
\begin{cxxdoc}

Maps the given fragmentlist fl into a vector of frames.

\end{cxxdoc}
\end{cxxfunction}
\begin{cxxfunction}
{AVTelos}
        {vt\_create}
        {(char*\ szTelos)}
        { Maps the given Telos text szTelos into a vector of frames.}
        {29}
\cxxParameter{
\begin{tabular}[t]{lp{0.5\textwidth}}
{\tt\strut szTelos} & should be a string of correct Telos
\end{tabular}}
\cxxReturn{
\begin{tabular}[t]{lp{0.5\textwidth}}
{\tt\strut NULL} & if the szTelos fails the parsing process else it contains the AVTelos with all its componends
\end{tabular}}
\begin{cxxdoc}

Maps the given Telos text szTelos into a vector of frames.

\end{cxxdoc}
\end{cxxfunction}
\begin{cxxfunction}
{void}
        {vt\_destroy}
        {(AVTelos\ avtFrames)}
        { Disposes the given vector tree.}
        {30}
\cxxParameter{
\begin{tabular}[t]{lp{0.5\textwidth}}
{\tt\strut avtFrames} & points to the vector tree structure
\end{tabular}}
\begin{cxxdoc}

Disposes the given vector tree.

\end{cxxdoc}
\end{cxxfunction}
\begin{cxxfunction}
{PReport}
        {rep\_create}
        {(PVTelos\ pvtFrame,\ char**\ aszCategories)}
        { Filters all attributes to those properties which belong to all given categories at the same time.}
        {31}
\cxxParameter{
\begin{tabular}[t]{lp{0.5\textwidth}}
{\tt\strut pvtFrame} & points to a single Frame, which must exists@pararm aszCategories should be a NULL terminated vector of the categories to filter as a conjunction.
\end{tabular}}
\cxxReturn{
\begin{tabular}[t]{lp{0.5\textwidth}}
{\tt\strut In} & each case a report will be created, even if the result is empty.
\end{tabular}}
\begin{cxxdoc}

Filters all attributes to those properties which belong to all given categories at the same time.
Note: If there is a category wrong typed, it has the effect that result will always be an empty vector with NULL at index 0.

\end{cxxdoc}
\end{cxxfunction}
\begin{cxxfunction}
{void}
        {rep\_destroy}
        {(PReport\ prepReport)}
        { Disposes the given report structure.}
        {32}
\cxxParameter{
\begin{tabular}[t]{lp{0.5\textwidth}}
{\tt\strut prepReport} & points to the report which should be disposed
\end{tabular}}
\begin{cxxdoc}

Disposes the given report structure. Should be called to free the result of rep\_create.

\end{cxxdoc}
\end{cxxfunction}
\begin{cxxfunction}
{char*}
        {getValueOfLabel}
        {(PVTelos\ pvtFrame,\ char*\ szLabel)}
        { A simple service routine, which support the access on frames.}
        {33}
\cxxParameter{
\begin{tabular}[t]{lp{0.5\textwidth}}
{\tt\strut pvtFrame} & points to a single frame, which should be analyzed
\\
{\tt\strut szLabel} & contains the keyword, which should be searched in the labels
\end{tabular}}
\cxxReturn{
\begin{tabular}[t]{lp{0.5\textwidth}}
{\tt\strut the} & string containing the value according to the given labelor NULL if no appropriate value was found
\end{tabular}}
\begin{cxxdoc}

A simple service routine, which support the access on frames.

\end{cxxdoc}
\end{cxxfunction}
\begin{cxxfunction}
{char**}
        {getCategories}
        {(PVTelos\ pvtFrame)}
        { Lists the categories as flat list, each elemant appears only once.}
        {34}
\cxxParameter{
\begin{tabular}[t]{lp{0.5\textwidth}}
{\tt\strut pvtFrame} & points to a single frame, which should be analyzed
\end{tabular}}
\cxxReturn{
\begin{tabular}[t]{lp{0.5\textwidth}}
{\tt\strut array} & of strings with the categories
\end{tabular}}
\begin{cxxdoc}

Lists the categories as flat list, each elemant appears only once.
The categories will be ordered by their appearance.

\end{cxxdoc}
\end{cxxfunction}
\begin{cxxfunction}
{void}
        {destroyASZ}
        {(char**\ asz)}
        { Disposes an asz (array of strings) structure.}
        {35}
\cxxParameter{
\begin{tabular}[t]{lp{0.5\textwidth}}
{\tt\strut asz} & the array of strings to be disposed
\end{tabular}}
\begin{cxxdoc}

Disposes an asz (array of strings) structure.

\end{cxxdoc}
\end{cxxfunction}



\section{te\_cursor.h}
\label{sec:cursorh}

\begin{cxxunion}
{struct\ }
        {te\_framecursor}
        {}
        { A cursor for a Telos frame}
        {36}
\begin{cxxnames}
\cxxitem{te\_SMLfragmentList*}
        {flAll}
        {}
        {Parsed telos frames in a fragment list }
        {36.1}
\cxxitem{te\_SMLfragmentList*}
        {flCur}
        {}
        {fragment list cursor }
        {36.2}
\cxxitem{te\_ClassList*}
        {clCurOmega}
        {}
        {cursor for inOmega }
        {36.3}
\cxxitem{te\_ClassList*}
        {clCurIn}
        {}
        {cursor for in }
        {36.4}
\cxxitem{te\_ClassList*}
        {clCurIsA}
        {}
        {cursor for isA }
        {36.5}
\cxxitem{AttrDeclList*}
        {alCur}
        {}
        {cursor for attribute declarations  }
        {36.6}
\cxxitem{AttrClassList*}
        {clCurCategory}
        {}
        {cursor for attribute categories (within one attribute declaration) }
        {36.7}
\cxxitem{PropertyList*}
        {plCur}
        {}
        {cursor for property list (within one attribute declaration) }
        {36.8}
\cxxitem{AttrDeclList*}
        {alChecked}
        {}
        {internal filter: tests if alCur was checked }
        {36.9}
\end{cxxnames}
\begin{cxxdoc}

A cursor for a Telos frame
\end{cxxdoc}
\begin{cxxvariable}
{te\_SMLfragmentList*}
        {flAll}
        {}
        {Parsed telos frames in a fragment list }
        {36.1}
\end{cxxvariable}
\begin{cxxvariable}
{te\_SMLfragmentList*}
        {flCur}
        {}
        {fragment list cursor }
        {36.2}
\end{cxxvariable}
\begin{cxxvariable}
{te\_ClassList*}
        {clCurOmega}
        {}
        {cursor for inOmega }
        {36.3}
\end{cxxvariable}
\begin{cxxvariable}
{te\_ClassList*}
        {clCurIn}
        {}
        {cursor for in }
        {36.4}
\end{cxxvariable}
\begin{cxxvariable}
{te\_ClassList*}
        {clCurIsA}
        {}
        {cursor for isA }
        {36.5}
\end{cxxvariable}
\begin{cxxvariable}
{AttrDeclList*}
        {alCur}
        {}
        {cursor for attribute declarations  }
        {36.6}
\end{cxxvariable}
\begin{cxxvariable}
{AttrClassList*}
        {clCurCategory}
        {}
        {cursor for attribute categories (within one attribute declaration) }
        {36.7}
\end{cxxvariable}
\begin{cxxvariable}
{PropertyList*}
        {plCur}
        {}
        {cursor for property list (within one attribute declaration) }
        {36.8}
\end{cxxvariable}
\begin{cxxvariable}
{AttrDeclList*}
        {alChecked}
        {}
        {internal filter: tests if alCur was checked }
        {36.9}
\end{cxxvariable}
\end{cxxunion}
\begin{cxxentry}
{typedef\ struct\ }
        {te\_framecursor}
        {TFrameCursor}
        {Type for te\_framecursor }
        {37}
\end{cxxentry}
\begin{cxxentry}
{typedef\ }
        {TFrameCursor*}
        {PFrameCursor}
        {Type for a pointer to TFrameCursor }
        {38}
\end{cxxentry}
\begin{cxxfunction}
{PFrameCursor}
        {te\_createCursor}
        {(\ te\_SMLfragmentList*\ fl)}
        { Creates and initializes a cursor structure for the given smlfragmentlist.}
        {39}
\begin{cxxdoc}

Creates and initializes a cursor structure for the given smlfragmentlist.
and returns a pointer to it.
\end{cxxdoc}
\end{cxxfunction}
\begin{cxxfunction}
{void}
        {te\_destroyCursor}
        {(PFrameCursor\ pfc)}
        { Deallocates a cursor structure}
        {40}
\begin{cxxdoc}

Deallocates a cursor structure
\end{cxxdoc}
\end{cxxfunction}
\begin{cxxfunction}
{void}
        {te\_resetFrame}
        {(PFrameCursor\ pfc)}
        { Sets the frame cursor to the first frame and resets all sub cursors}
        {41}
\begin{cxxdoc}

Sets the frame cursor to the first frame and resets all sub cursors
\end{cxxdoc}
\end{cxxfunction}
\begin{cxxfunction}
{int}
        {te\_nextFrame}
        {(PFrameCursor\ pfc)}
        { Sets the frame cursor to the next frame in the list and resets the Omega, IsA, In, AttrDecl, Category and Property cursors.}
        {42}
\cxxReturn{
\begin{tabular}[t]{lp{0.5\textwidth}}
{\tt\strut true} & (non-zero) if there is a next element
\end{tabular}}
\begin{cxxdoc}

Sets the frame cursor to the next frame in the list and resets the
Omega, IsA, In, AttrDecl, Category and Property cursors.

\end{cxxdoc}
\end{cxxfunction}
\begin{cxxfunction}
{char*}
        {te\_retOID}
        {(PFrameCursor\ pfc)}
        { Returns the OID of a frame as plain string, even if it is a select expression.}
        {43}
\begin{cxxdoc}

Returns the OID of a frame as plain string, even if it is
a select expression. Memory for this string has to be deallocated
by the caller.
\end{cxxdoc}
\end{cxxfunction}
\begin{cxxfunction}
{void}
        {te\_resetOmega}
        {(PFrameCursor\ pfc)}
        { Resets the Omega cursor}
        {44}
\begin{cxxdoc}

Resets the Omega cursor
\end{cxxdoc}
\end{cxxfunction}
\begin{cxxfunction}
{int}
        {te\_nextOmega}
        {(PFrameCursor\ pfc)}
        { Sets the omega cursor to the next element.}
        {45}
\cxxReturn{
\begin{tabular}[t]{lp{0.5\textwidth}}
{\tt\strut true} & (non-zero) if there is a next element
\end{tabular}}
\begin{cxxdoc}

Sets the omega cursor to the next element.

\end{cxxdoc}
\end{cxxfunction}
\begin{cxxfunction}
{char*}
        {te\_retOmega}
        {(PFrameCursor\ pfc)}
        { Returns the omega object as string }
        {46}
\begin{cxxdoc}

Returns the omega object as string

\end{cxxdoc}
\end{cxxfunction}
\begin{cxxfunction}
{void}
        {te\_resetIsA}
        {(PFrameCursor\ pfc)}
        { Resets the Isa cursor}
        {47}
\begin{cxxdoc}

Resets the Isa cursor
\end{cxxdoc}
\end{cxxfunction}
\begin{cxxfunction}
{int}
        {te\_nextIsA}
        {(PFrameCursor\ pfc)}
        { Sets the IsA cursor to the next element.}
        {48}
\cxxReturn{
\begin{tabular}[t]{lp{0.5\textwidth}}
{\tt\strut true} & (non-zero) if there is a next element
\end{tabular}}
\begin{cxxdoc}

Sets the IsA cursor to the next element.

\end{cxxdoc}
\end{cxxfunction}
\begin{cxxfunction}
{char*}
        {te\_retIsA}
        {(PFrameCursor\ pfc)}
        { Returns the IsA object as string }
        {49}
\begin{cxxdoc}

Returns the IsA object as string

\end{cxxdoc}
\end{cxxfunction}
\begin{cxxfunction}
{void}
        {te\_resetIn}
        {(PFrameCursor\ pfc)}
        { Resets the In cursor}
        {50}
\begin{cxxdoc}

Resets the In cursor
\end{cxxdoc}
\end{cxxfunction}
\begin{cxxfunction}
{int}
        {te\_nextIn}
        {(PFrameCursor\ pfc)}
        { Sets the In cursor to the next element.}
        {51}
\cxxReturn{
\begin{tabular}[t]{lp{0.5\textwidth}}
{\tt\strut true} & (non-zero) if there is a next element
\end{tabular}}
\begin{cxxdoc}

Sets the In cursor to the next element.

\end{cxxdoc}
\end{cxxfunction}
\begin{cxxfunction}
{char*}
        {te\_retIn}
        {(PFrameCursor\ pfc)}
        { Returns the In object as string }
        {52}
\begin{cxxdoc}

Returns the In object as string

\end{cxxdoc}
\end{cxxfunction}
\begin{cxxfunction}
{void}
        {te\_resetAttrDecl}
        {(PFrameCursor\ pfc)}
        { Sets the attr decl block cursor to the first attr decl block  in the current frame and resets the sub-cursors Category and Property}
        {53}
\begin{cxxdoc}

Sets the attr decl block cursor to the first attr decl block  in the
current frame and resets the sub-cursors Category and Property
\end{cxxdoc}
\end{cxxfunction}
\begin{cxxfunction}
{int}
        {te\_nextAttrDecl}
        {(PFrameCursor\ pfc)}
        { Sets the Property cursor to the next Property class in the attr decl block frame and resets the Category and Property cursors.}
        {54}
\cxxReturn{
\begin{tabular}[t]{lp{0.5\textwidth}}
{\tt\strut true} & (non-zero) if there is a next element
\end{tabular}}
\begin{cxxdoc}

Sets the Property cursor to the next Property class in the attr decl
block frame and resets the Category and Property cursors.

\end{cxxdoc}
\end{cxxfunction}
\begin{cxxfunction}
{void}
        {te\_resetCategory}
        {(PFrameCursor\ pfc)}
        { Resets the category cursor}
        {55}
\begin{cxxdoc}

Resets the category cursor
\end{cxxdoc}
\end{cxxfunction}
\begin{cxxfunction}
{int}
        {te\_nextCategory}
        {(PFrameCursor\ pfc)}
        { Sets the category cursor to the next element.}
        {56}
\cxxReturn{
\begin{tabular}[t]{lp{0.5\textwidth}}
{\tt\strut true} & (non-zero) if there is a next element
\end{tabular}}
\begin{cxxdoc}

Sets the category cursor to the next element.

\end{cxxdoc}
\end{cxxfunction}
\begin{cxxfunction}
{char*}
        {te\_retCategory}
        {(PFrameCursor\ pfc)}
        { Returns the category as string }
        {57}
\begin{cxxdoc}

Returns the category as string

\end{cxxdoc}
\end{cxxfunction}
\begin{cxxfunction}
{void}
        {te\_resetProperty}
        {(PFrameCursor\ pfc)}
        { Resets the property cursor}
        {58}
\begin{cxxdoc}

Resets the property cursor
\end{cxxdoc}
\end{cxxfunction}
\begin{cxxfunction}
{int}
        {te\_nextProperty}
        {(PFrameCursor\ pfc)}
        { Sets the category cursor to the next element.}
        {59}
\cxxReturn{
\begin{tabular}[t]{lp{0.5\textwidth}}
{\tt\strut true} & (non-zero) if there is a next element
\end{tabular}}
\begin{cxxdoc}

Sets the category cursor to the next element.

\end{cxxdoc}
\end{cxxfunction}
\begin{cxxfunction}
{int}
        {te\_filterPropertyByCategory}
        {(PFrameCursor\ pfc,\ char*\ category)}
        { Lists all properties that are of the type "category".}
        {60}
\begin{cxxdoc}

Lists all properties that are of the type "category".
The usage of this function is similar to the function
te\_filterPropertyByCategories. The only diffrence is the simplicity of
the second parameter for the case that you only need to filter with
one category.
\end{cxxdoc}
\end{cxxfunction}
\begin{cxxfunction}
{int}
        {te\_filterPropertyByCategories}
        {(PFrameCursor\ pfc,\ char*\ categories[])}
        { Lists all properties that matches all types of categories If the categories are empty then any category matches.}
        {61}
\begin{cxxdoc}

Lists all properties that matches all types of categories
If the categories are empty then any category matches.
In difference to the nextXXX functions, these function should be
called before the first access via te\_retLabel or te\_retValue,
because it must search the first valid AttrDecl.
This means at the beginning you should call:
"te\_resetAttrDecl( ... );"          AND
"te\_filterPropertyByCategories( ... );"
\end{cxxdoc}
\end{cxxfunction}
\begin{cxxfunction}
{char*}
        {te\_filterPropertyByLabel}
        {(PFrameCursor\ pfc,\ char*\ )}
        { Lists the value of the property with label of the current frame.}
        {62}
\begin{cxxdoc}

Lists the value of the property with label of the current frame.
This results only one value which is a new created string.
Note that the caller has to dispose the return value !
\end{cxxdoc}
\end{cxxfunction}
\begin{cxxfunction}
{char*}
        {te\_retLabel}
        {(PFrameCursor\ pfc)}
        { Returns the current label of the current property in the current decl block or NULL}
        {63}
\begin{cxxdoc}

Returns the current label of the current property in the current
decl block or NULL
\end{cxxdoc}
\end{cxxfunction}
\begin{cxxfunction}
{char*}
        {te\_retValue}
        {(PFrameCursor\ pfc)}
        { Returns the current value of the current property in the current decl block or NULL}
        {64}
\begin{cxxdoc}

Returns the current value of the current property in the current
decl block or NULL
\end{cxxdoc}
\end{cxxfunction}

